\subsection{Evaluation of the finite inequalities}
\label{sec:window-finite-evaluation}

The finite inequalities in Corollary~\ref{cor:window-conditional}
hold for both rows of its table.  We give the polynomial construction,
the arithmetic error bounds and the exact positivity calculation.  The
matrix to be tested retains the full residual
$J-C^*G_K^{-1}C$ throughout.

\subsubsection{Finite polynomial construction}

Fix either $T=3/4$ or $T=81/100$, and put $L=2T$.  The
shifted Legendre polynomial has the integer coefficients
\begin{equation}
 p_j(t)=\sum_{k=0}^j b_{kj}t^k,\qquad
 b_{kj}=(-1)^{j-k}\binom jk\binom{j+k}k.
 \label{eq:window-computation-legendre}
\end{equation}
The scalar in the approximating operator is
$c_*=-\gamma_* -\log(2\pi T)$, where $\gamma_*$ is the finite
decimal evaluation of $C_{1000,40}-\log1000$ specified below.
It is an exact rational scalar once chosen.  Its distance to
$\gamma$ is counted in the bounded operator error, separately
from errors in evaluating the moments of the chosen operator.

Let $\mathcal R_{T,64}$ be the integral operator with kernel
$\sum_{k=0}^{64}r_k|x-y|^k$, using
\eqref{eq:window-r-recurrence}.  For a monomial in the coordinate
$t=(x+T)/L$, its image is the following polynomial:
\begin{equation}
 \mathcal R_{T,64}(t^m)=
 \sum_{k=0}^{64}r_kL^{k+1}
 \left\{
 \sum_{\ell=0}^k
  \frac{(-1)^{k-\ell}\binom k\ell}{m+\ell+1}
       t^{k-\ell}
 +\mathbf1_{\{k\text{ odd}\}}
       \frac{2m!k!}{(m+k+1)!}t^{m+k+1}
 \right\}.
 \label{eq:window-computation-kernel}
\end{equation}
To check this formula, first integrate $(s-t)^ks^m$ over
$0<s<1$.  When $k$ is odd, the correction on $s<t$ is
$2(t-s)^ks^m$; its integral is the displayed beta integral.
For even $k$, no correction is needed.

Put
\[
 v_\pm(t)=\sum_{k=0}^{32}\frac{(\pm(Lt-T)/2)^k}{k!}.
\]
The polar approximation maps $p_j$ to
\[
 v_+(t)L\int_0^1v_-(s)p_j(s)\,ds
 +v_-(t)L\int_0^1v_+(s)p_j(s)\,ds.
\]
These coefficients are rational.  For every retained prime power,
write $a_n=\log n/L$.  Its two polynomial pieces are
\[
 -c_n\mathbf1_{\{t<1-a_n\}}p_j(t+a_n)
 -c_n\mathbf1_{\{t>a_n\}}p_j(t-a_n).
\]
In particular $c_4=\log2/2$; the $n=5$ terms occur in the
second window.  Adding these pieces to
$(H_j+c_*)p_j-\mathcal R_{T,64}p_j$ and the polar image
gives the polynomial $q_j$ in
\eqref{eq:window-piecewise-column} on every interval.
All translation endpoints and all mixed products are retained.

The matrices are obtained by integrating
\[
 M_{ij}=L\int_0^1p_i(Vp_j+q_j)\,dt,\qquad
 C_{nj}=L\int_0^1p_n(Vp_j+q_j)\,dt,
\]
\begin{equation}
 J_{ij}=L\int_0^1
 \bigl(V^2p_ip_j+Vp_iq_j+Vp_jq_i+q_iq_j\bigr)\,dt.
 \label{eq:window-computation-gram}
\end{equation}
The $V^2$ integral has the full interval as its domain.
The other terms are split at the translation endpoints.
Ordinary monomial integrals and
\eqref{eq:window-log-moments} therefore determine every entry.
Here $M,J$ are real symmetric and $C$ is real; the resulting
matrices represent the original Hermitian forms on complex
coefficient vectors.
No numerical quadrature is used.

\subsubsection{Arithmetic enclosures}

Use $p=520$ significant decimal digits for centers, with rounding
to nearest and ties to an even final digit.  Products and sums
of nonnegative error bounds are rounded upward at the same
precision.  Polynomial indices are summed in increasing order,
and intervals are added from left to right.  The exact rational
coefficient formulas precede their conversion to decimals.

Use the identity
\[
 \pi=16\arctan(1/5)-4\arctan(1/239),
\]
which follows by tangent addition with the angles in the first
quadrant.  For each arctangent use
\[
 \arctan z=\sum_{k=0}^{499}\frac{(-1)^kz^{2k+1}}{2k+1}+\eta,
 \qquad |\eta|\le\frac{z^{1001}}{1001}.
\]
The omitted contribution to $\pi$ is less than $10^{-699}$.
After multiplying a positive logarithm argument
by a power of two to put it in $[1,2)$, use
\[
 \log y=2\sum_{j=0}^{599}\frac{z^{2j+1}}{2j+1}+\rho,
 \qquad z=\frac{y-1}{y+1},\qquad
 |\rho|\le\frac{2(1/3)^{1201}}{1201(1-1/9)}<10^{-572}.
\]
For $\log2$, take $z=1/3$ directly.  Restore the integer
multiple of this value after reducing any other argument.
Each scalar evaluation has fewer than $10^4$ arithmetic operations;
the reduced denominators are bounded away from zero.  The required
logarithm arguments lie between $0.006$ and $1000$.
Bounding their rounding errors by decimal units gives an error
less than $10^{-500}$.  The enlarged common radius
\[
 r_c=10^{-490}
\]
covers their finite combinations.  Square-root centers are enclosed
with radius $10^{-500}$ by squaring both rational endpoints.
For $\gamma_*$, evaluate the exact rational $C_{1000,40}$ and
the preceding logarithm formula at this precision.  Equation
\eqref{eq:window-gamma-error} gives
\[
 |\gamma-\gamma_*|\le E_{1000,40}+10^{-490}
                  <2.502\cdot10^{-187}.
\]

The proper breakpoints and their complements exceed $0.07$ in
the first window and $0.006$ in the second.  Their errors are at
most $2r_c$.  An ordinary primitive
$t^{k+1}/(k+1)$ has derivative of absolute value at most one
on $[0,1]$, so endpoint perturbations contribute at most $4r_c$.
The radius $10^{-480}$ also covers center rounding for every
ordinary moment used here.

For a logarithmic primitive
\[
 F_k(t)=t^{k+1}\left(\frac{\log t}{k+1}
                              -\frac1{(k+1)^2}\right),
 \qquad F_k'(t)=t^k\log t,
\]
the value at zero is the exact limit.  At every proper endpoint,
$|\log t|<6$ and $1/t<170$.  Its input and rounding errors
are less than $10^{-487}$.  The substitution $t\mapsto1-t$
has absolute binomial coefficient sum $2^k$.  The maximal
degrees are $191$ and $319$, with
\[
 2^{191}<10^{58},\qquad 2^{319}<10^{97}.
\]
Rounding the alternating sum is bounded by
$2(k+1)10^{\lceil k\log_{10}2\rceil+4-520}$.
Consequently radii $10^{-420}$ and $10^{-380}$ cover the two
mixed-logarithm calculations, including cancellation.
For the global squared-logarithm moments, $H_j<7$,
$H_j^{(2)}<2$, and the remaining scalar factors are bounded by
$100$; radius $10^{-480}$ suffices.

The translation coefficient matrices use the power expansion of
$p_j(t\pm a_n)$.  Since $|a_n|<1$ and
$|(a_n^r)'|\le r$, the bounds
\[
 64\,2^{63}<10^{21},\qquad
 128\,2^{127}<10^{41}
\]
justify their entry radii $10^{25}r_c$ and $10^{45}r_c$.
The diagonal polynomial coefficients use radius
$\max|b_{kj}|r_c+10^{-450}$ in the first window and
$\max|b_{kj}|r_c+10^{-380}$ in the second.
These bounds include center arithmetic as well as input errors.

Here are the propagation rules.  If centers $\widehat A,
\widehat B$ have uniform entry radii $r_A,r_B$, their product
of inner dimension $k$ has radius
\begin{equation}
 r_A\|\widehat B\|_1+r_B\|\widehat A\|_\infty
       +kr_Ar_B+r_{\mathrm{round}}.
 \label{eq:window-computation-radius}
\end{equation}
The two matrix norms in this formula are the maximum absolute
column and row sums.  If the center magnitudes are less than
$10^{e_A}$ and $10^{e_B}$, a valid rounding term is
\[
 r_{\mathrm{round}}=
 2k\,10^{e_A+e_B+d(k)+3-p},\qquad
 d(k)=1+\lfloor\log_{10}k\rfloor.
\]
It bounds every rounded product and partial sum.  Addition uses
$r_A+r_B$ and a rounding allowance
$10^{\max(e_A,e_B)+4-p}$.  Multiplication by a scalar center
$\widehat c$ with radius $r_c'$ uses
\[
 |\widehat c|r_A+\max|\widehat A_{ij}|r_c'
                  +r_Ar_c'+r_{\mathrm{round}}.
\]
Here one may take $r_{\mathrm{round}}=10^{e_A+e_c+3-p}$
when $|\widehat c|<10^{e_c}$.  Conversion of a rational
coefficient of magnitude less than $10^e$ uses radius $10^{e+3-p}$.
For separate row weights, use the largest absolute weight in
this bound.  These formulas do not require independence of the
input errors.

The matrix assembly in \eqref{eq:window-computation-gram}, with
these upward error rules, gives the following strict rational
upper bounds before shortening any centers:
\[
 \begin{array}{c|c|c|c}
 T&\text{intervals}&r_M=r_C&r_J\\\hline
 3/4&7&10^{-223}&10^{-312}\\
 81/100&9&10^{-154}&2\cdot10^{-159}
 \end{array}
\]
Symmetric entries are averaged and opposite-parity entries are
set to their exact value zero.  For $T=81/100$, rounding the
centers to $90$ significant digits adds less than $10^{-88}$
per entry, since every center has magnitude less than $100$.
Adding $10^{-85}$ to the preceding radii gives radii less than
$2\cdot10^{-85}$.  In both windows the resulting enclosures
therefore satisfy the $10^{-80}$ entry requirement used in
\eqref{eq:window-rational-error}.

The scalar estimates give, with the complete arithmetic lists,
\[
 a(3/4)<4.30094,\qquad a(81/100)<5.16201132.
\]
The kernel recurrence through degree $128$, the polar remainder
and the preceding Euler-constant enclosure give total bounded
operator errors less than $4\cdot10^{-22}$ and
$7.3\cdot10^{-20}$, respectively.  The exact harmonic sums
then give $d_{64}>0.44$ and $d_{128}>0.27$ at the targets in
Corollary~\ref{cor:window-conditional}.

\subsubsection{Exact rational positivity}

Round the enclosed centers downward to the grid $10^{-80}$
and the inverse weights downward to the grid $10^{-100}$ as
in \eqref{eq:window-rational-error}.  For both windows its exact
rational error bound satisfies
\[
 E_F<10^{-68}<\kappa:=10^{-60}.
\]
It remains to test the rational matrix $\widehat F-\kappa I$.
The two parity blocks have orders $32$ and $32$, or $64$ and
$64$, respectively.

Here is a fraction-free form of the test.  For a parity block
$A$, let $D$ be the least common multiple of its entry
denominators and set $B^{(0)}=DA$.  Put $p_{-1}=1$.
At step $k$, set $p_k=B^{(k)}_{kk}$ and, for $i,j>k$, put
\begin{equation}
 B^{(k+1)}_{ij}=
 \frac{p_kB^{(k)}_{ij}-B^{(k)}_{ik}B^{(k)}_{kj}}{p_{k-1}}.
 \label{eq:window-computation-bareiss}
\end{equation}
The determinant identity for adjacent minors proves that each
division is exact and that $p_k$ is the leading principal
minor of order $k+1$ of $DA$.  Thus $p_k/(Dp_{k-1})$ is
the corresponding rational $LDL^*$ pivot of $A$.

Applying \eqref{eq:window-computation-bareiss} to the rational
matrices just specified gives strictly positive pivots throughout.
Simple strict lower bounds for every pivot are
\[
 \begin{array}{c|c|c|c}
 T&\lambda&\text{even block}&\text{odd block}\\\hline
 3/4&0&9\cdot10^{-5}&2\cdot10^{-4}\\
 81/100&10^{-18}&3\cdot10^{-5}&10^{-4}
 \end{array}
\]
These are pivots in the unnormalized Legendre coefficient
matrix, not spectral lower endpoints in the physical norm.
The latter are the targets $\lambda$ already incorporated in
$\widehat F$.  Positivity of all leading principal minors
proves $\widehat F-\kappa I>0$ in both cases.

\begin{theorem}[Two complete finite-window bounds]
\label{thm:window-certified}
On the original logarithmic form domains,
\[
 Q_{3/4}(g)\ge0\quad(g\in\mathcal Q_{3/4}),\qquad
 Q_{81/100}(g)\ge10^{-18}\|g\|_2^2
                \quad(g\in\mathcal Q_{81/100}).
\]
Every prime power, the signed polar cross term and the complete
remote Gram matrix are included in these inequalities.
\end{theorem}
\begin{proof}
The polynomial construction and propagated radii enclose the
entries of the specified approximating operator.  The scalar
estimates bound its distance from the original operator and
give positive tail denominators.  The exact rational calculation
proves the finite positivity hypotheses of
Corollary~\ref{cor:window-conditional}, including its rounding
cushion.  That corollary yields the two inequalities on the
entire original form domains.
\end{proof}

This calculation treats two fixed support windows.  It supplies
neither an estimate on arbitrarily large supports nor the cofinal
signed-matrix inequality.
